

% ============================================================
%  abel.tex  —  FRAGMENT, included via \input from main.tex
%  Story: how the MDP changed from v1 (mdp_old) to v2 (mdp)
%  Uses colors/styles defined in main.tex.
% ============================================================

% ============================================================
\section{MDP Formulation}
% ============================================================

% ---- Slide 0: Define the MDP -------------------------------
\begin{frame}{Defining the Problem as an MDP}
  \small
  A cluster of $N$ identical nodes runs a stream of arriving jobs. Each job asks
  for a number of nodes and a wall-clock runtime. At every decision point the
  scheduler may decide \emph{which job runs next} from a set of $K$ visible jobs.

  \vspace{0.3em}
  \begin{block}{Why a Markov Decision Process?}
    \small
    Scheduling is \textbf{sequential}: each choice changes the cluster the next
    one is made on, and its payoff (utilisation, delay) shows up later. An MDP
    captures exactly this --- an agent observes a \emph{state}, takes an
    \emph{action}, gets a \emph{reward}, and moves to the next state:
    $$ \mathcal{M} = (\underbrace{\mathcal{S}}_{\text{state}},\ \underbrace{\mathcal{A}}_{\text{action}},\ \underbrace{P}_{\text{transition}},\ \underbrace{R}_{\text{reward}},\ \underbrace{\gamma}_{\text{discount}}) $$
  \end{block}

  \vspace{0.1em}
  We learn a policy $\pi_\theta(a_t\mid s_t)$ with RL that maximises long-run
  reward. The next slides walk through each component
  $(\mathcal{S},\mathcal{A},P,R)$ in turn.
\end{frame}

% ---- Slide 1: Roadmap of what changed ----------------------
\begin{frame}{Walking Through the MDP: What Changed}
  \small
  We keep the \textbf{action} and \textbf{dynamics}, and redesign the
  \textbf{state} and \textbf{reward} to be \alert{bounded and load-agnostic}.

  \vspace{0.4em}
  \begin{center}
  \renewcommand{\arraystretch}{1.25}
  \begin{tabular}{l c c}
    \toprule
    \textbf{Component} & \textbf{Earlier (v1)} & \textbf{Now (v2)} \\
    \midrule
    State      & $12$-dim aggregated, raw scales   & \textcolor{SoftCoral}{\textbf{$2K{+}2$ per-job, normalised}} \\
    Action     & Discrete$(K{+}1)$ promote         & unchanged \\
    Transition & $\Delta t$ step, stochastic arrivals & unchanged \\
    Reward     & $w_u u_t - w_Q Q_t$ (unbounded)   & \textcolor{SoftCoral}{\textbf{$-w\,\tfrac{\min(D_t,n_{\text{free}})}{N}\in[-1,0]$}} \\
    Horizon    & truncate at $T$ or drained        & unchanged \\
    \bottomrule
  \end{tabular}
  \end{center}

  \vspace{0.3em}
  \centering
  \small\textcolor{DeepBlue}{Comparison slides follow \emph{only} for the two components that changed.}
\end{frame}

% ---- Slide 2: State change (CHANGED) -----------------------
\begin{frame}{State $\;s_t$ \hfill \small\textcolor{SoftCoral}{(redesigned)}}
  \begin{columns}[T]
    \column{0.5\textwidth}
      \begin{block}{Earlier (v1): aggregated summary}
        \small
        $s_t \in \mathbb{R}^{12}$, \textbf{raw} unbounded values.
        \begin{itemize}
          \item Queue collapsed to \emph{means}: $\bar{q}_t,\ \bar{w}_t$
          \item Native scales (e.g. elapsed time in hours)
          \item $\mathcal{S} = \text{Box}\big([-1,10^{6}]^{12}\big)$
        \end{itemize}
      \end{block}

    \column{0.5\textwidth}
      \begin{block}{Now (v2): per-job, normalised}
        \small
        $s_t \in \mathbb{R}^{2K+2}$, all scaled to $[0,1]$:
        $$
          s_t = \Big[\tfrac{t}{T},\ \tfrac{n_{\text{free}}}{N},\
          \big(\tfrac{q_i}{N}\big)_{i=1}^{K},\
          \big(\tfrac{w_i}{T}\big)_{i=1}^{K}\Big]
        $$
        \begin{itemize}
          \item Exposes \emph{each} visible job, not just a mean
          \item Empty slots padded with $-1$ sentinel
          \item $\mathcal{S} = \text{Box}\big([0,1]^2\times[-1,1]^{2K}\big)$
        \end{itemize}
      \end{block}
  \end{columns}

  \vspace{0.2em}
  \centering
  \small\textcolor{DeepBlue}{\textbf{Why:}} bounded, stationary inputs $\Rightarrow$ stable learning and better generalisation.
\end{frame}

% ---- Slide 3: Action (UNCHANGED) ---------------------------
\begin{frame}{Action $\;a_t$ \hfill \small\textcolor{SkyBlue!70!black}{(unchanged)}}
  \small
  A single \textbf{discrete} choice over the visible queue:
  $$ \mathcal{A} = \{0, 1, 2, \dots, K\} = \text{Discrete}(K{+}1) $$
  \begin{itemize}
    \item $a_t = 0$ --- \textbf{no-op}: leave the queue ordering unchanged.
    \item $a_t = k \in \{1,\dots,K\}$ --- \textbf{promote} the $k$-th visible queued
          job to the front, so it is considered first on the next dispatch.
  \end{itemize}
  \begin{block}{Key idea}
    The agent never dispatches jobs directly --- it only \emph{reprioritises}.
    The underlying scheduler then greedily packs jobs that fit onto free nodes.
  \end{block}
\end{frame}

% ---- Slide 4: Transition (UNCHANGED) -----------------------
\begin{frame}{Transition $\;P(s_{t+1}\mid s_t, a_t)$ \hfill \small\textcolor{SkyBlue!70!black}{(unchanged)}}
  \small
  Each step advances the simulation by $\Delta t$ seconds:
  \begin{enumerate}
    \item \textbf{Apply the action} --- if $a_t = k > 0$, move that job to the head of the queue.
    \item \textbf{Step the scheduler} by $\Delta t$: dispatch queued jobs onto free
          nodes (greedy packing), advance running jobs, release nodes from jobs that finish.
    \item \textbf{New arrivals} --- jobs arrive stochastically (Poisson-like
          inter-arrival times set by the \emph{offered load}).
  \end{enumerate}
  \begin{block}{Source of stochasticity}
    Job arrivals and the workload sampled each episode make the next state only
    \emph{partially} predictable from $(s_t, a_t)$.
  \end{block}
\end{frame}

% ---- Slide 5: Reward change (CHANGED) ----------------------
\begin{frame}{Reward $\;R(s_t, a_t)$ \hfill \small\textcolor{SoftCoral}{(redesigned)}}
  \footnotesize
  \begin{columns}[T,onlytextwidth]
    \column{0.315\textwidth}
      \begin{block}{v1: util.\ $-$ queue}
        \begin{minipage}[t][3.5cm]{\linewidth}
        $$ r_t = w_u u_t - w_Q Q_t $$
        \begin{itemize}\setlength\itemsep{2pt}
          \item Rewards busy nodes, penalises backlog
          \item \alert{Unbounded} --- uses raw queue count $Q_t$
        \end{itemize}
        \end{minipage}
      \end{block}

    \hfill
    \column{0.315\textwidth}
      \begin{block}{v2: schedulable-idle}
        \begin{minipage}[t][3.5cm]{\linewidth}
        $$ r_t = -w\,\frac{\min(D_t, n_{\text{free}})}{N} $$
        \begin{itemize}\setlength\itemsep{2pt}
          \item Charges only for idle nodes the queue \emph{could} use
          \item Bounded $[-1,0]$, \textcolor{DeepBlue}{load-agnostic}
        \end{itemize}
        \end{minipage}
      \end{block}

    \hfill
    \column{0.315\textwidth}
      \begin{block}{v2: priority slowdown}
        \begin{minipage}[t][3.5cm]{\linewidth}
        $$ r^{\text{p}}_t = -w_p\,\widehat{\textstyle\sum_{j} \tfrac{p_j}{T_j}} $$
        \begin{itemize}\setlength\itemsep{2pt}
          \item Over \textbf{running \& queued} jobs ($p_j$=priority, $T_j$=runtime)
          \item Ranks by ratio $p_j/T_j$ --- \emph{priority-weighted shortest-job-first}
          \item Big job wins \textbf{only if} important enough: $p_j/T_j$ beats rivals
        \end{itemize}
        \end{minipage}
      \end{block}
  \end{columns}

  \vspace{0.4em}
  \centering
  \textcolor{DeepBlue}{\textbf{Why:}} bounded, well-scaled signals that target head-of-line blocking and priority directly.
\end{frame}

% ---- Slide 6: Horizon & Objective (UNCHANGED) --------------
\begin{frame}{Horizon \& Objective \hfill \small\textcolor{SkyBlue!70!black}{(unchanged)}}
  \small
  An episode \textbf{truncates} when either
  $$ t \geq T \quad(\text{24\,h elapsed}) \qquad\text{or}\qquad \text{queue \& running sets are empty (drained).} $$
  There is no absorbing failure state, so \texttt{terminated} is always \texttt{False};
  episodes end only by truncation.

  \begin{block}{Learning objective}
    The agent learns a policy $\pi_\theta(a_t \mid s_t)$ maximising the expected
    discounted return:
    $$ \max_{\theta}\; \mathbb{E}_{\pi_\theta}\!\left[\sum_{t=0}^{H-1} \gamma^{\,t}\, r_t\right] $$
    i.e.\ \textbf{keep the cluster busy} while draining the queue.
  \end{block}
\end{frame}

% ---- Slide 7: Why it matters -------------------------------
\begin{frame}{Why These Changes Matter}
  \begin{columns}[T]
    \column{0.52\textwidth}
      \textbf{From v1 to v2 we gain:}
      \begin{itemize}
        \item \textbf{Bounded signals} --- state in $[0,1]$, reward in $[-1,0]$
        \item \textbf{Load invariance} --- no dependence on raw queue size or job distribution
        \item \textbf{Richer observation} --- per-job features instead of a single mean
        \item \textbf{Direct objective} --- penalise \emph{schedulable} idle capacity (head-of-line blocking)
      \end{itemize}

    \column{0.48\textwidth}
      \begin{alertblock}{Net effect}
        A policy that learns to \textbf{promote jobs that fit} the free nodes,
        draining the queue faster --- and that \textbf{transfers} across
        workloads without re-tuning reward weights.
      \end{alertblock}
  \end{columns}
\end{frame}